% Options for packages loaded elsewhere
\PassOptionsToPackage{unicode}{hyperref}
\PassOptionsToPackage{hyphens}{url}
\documentclass[
  ignorenonframetext,
]{beamer}
\newif\ifbibliography
\usepackage{pgfpages}
\setbeamertemplate{caption}[numbered]
\setbeamertemplate{caption label separator}{: }
\setbeamercolor{caption name}{fg=normal text.fg}
\beamertemplatenavigationsymbolsempty
% remove section numbering
\setbeamertemplate{part page}{
  \centering
  \begin{beamercolorbox}[sep=16pt,center]{part title}
    \usebeamerfont{part title}\insertpart\par
  \end{beamercolorbox}
}
\setbeamertemplate{section page}{
  \centering
  \begin{beamercolorbox}[sep=12pt,center]{section title}
    \usebeamerfont{section title}\insertsection\par
  \end{beamercolorbox}
}
\setbeamertemplate{subsection page}{
  \centering
  \begin{beamercolorbox}[sep=8pt,center]{subsection title}
    \usebeamerfont{subsection title}\insertsubsection\par
  \end{beamercolorbox}
}
% Prevent slide breaks in the middle of a paragraph
\widowpenalties 1 10000
\raggedbottom
\AtBeginPart{
  \frame{\partpage}
}
\AtBeginSection{
  \ifbibliography
  \else
    \frame{\sectionpage}
  \fi
}
\AtBeginSubsection{
  \frame{\subsectionpage}
}
\usepackage{iftex}
\ifPDFTeX
  \usepackage[T1]{fontenc}
  \usepackage[utf8]{inputenc}
  \usepackage{textcomp} % provide euro and other symbols
\else % if luatex or xetex
  \usepackage{unicode-math} % this also loads fontspec
  \defaultfontfeatures{Scale=MatchLowercase}
  \defaultfontfeatures[\rmfamily]{Ligatures=TeX,Scale=1}
\fi
\usepackage{lmodern}
\ifPDFTeX\else
  % xetex/luatex font selection
\fi
% Use upquote if available, for straight quotes in verbatim environments
\IfFileExists{upquote.sty}{\usepackage{upquote}}{}
\IfFileExists{microtype.sty}{% use microtype if available
  \usepackage[]{microtype}
  \UseMicrotypeSet[protrusion]{basicmath} % disable protrusion for tt fonts
}{}
\makeatletter
\@ifundefined{KOMAClassName}{% if non-KOMA class
  \IfFileExists{parskip.sty}{%
    \usepackage{parskip}
  }{% else
    \setlength{\parindent}{0pt}
    \setlength{\parskip}{6pt plus 2pt minus 1pt}}
}{% if KOMA class
  \KOMAoptions{parskip=half}}
\makeatother
\usepackage{longtable,booktabs,array}
\usepackage{caption}
\captionsetup[table]{skip=6pt}
\newcounter{none} % for unnumbered tables
\usepackage{calc} % for calculating minipage widths
\usepackage{caption}
% Make caption package work with longtable
\makeatletter
\def\fnum@table{\tablename~\thetable}
\makeatother
\setlength{\emergencystretch}{3em} % prevent overfull lines
\providecommand{\tightlist}{%
  \setlength{\itemsep}{0pt}\setlength{\parskip}{0pt}}
% Auto-generated by infrastructure/rendering/_pdf_latex_helpers.py::extract_math_font_preamble.
% Minimal header passed to Pandoc via -H so Beamer slides pick up the same math font
% as the combined-PDF path. Adjust the manuscript's preamble.md to change the font globally.
\usepackage{unicode-math}
\setmathfont{latinmodern-math.otf}

% Manuscript macro/package fallbacks (newcommand -> providecommand).
% Core mathematics
\usepackage{amsmath}
\usepackage{amssymb}
% Document layout
% Tables
\usepackage{booktabs}
% Algorithm typesetting
% Code listings
\usepackage{listings}
% Typography and formatting
% Cross-references and citations
% ── Unicode-capable mono font for code listings ──────────────────────
% JuliaMono (TeX Live 2026) covers the full math/Greek glyph set used in
% scientific code blocks; the default lmmono lacks \alpha/\mu/\partial/\nabla.
%
% The hardcoded TeX Live 2026 path below is portable to most macOS / Linux
% TeX Live installs that placed JuliaMono under the default texmf-dist path.
% If your TeX Live version differs (e.g., 2025, 2027) or your distribution
% installs fonts elsewhere, comment out the explicit Path = line — fontspec
% will then search the system font path. If JuliaMono is not installed at
% all, comment out the whole block; xelatex falls back to lmmono with
% reduced Unicode coverage but the document still compiles.
% Math font for unicode-math: Latin Modern Math (TeX Live) has full BMP
% coverage including U+2223 (\mid), U+226A/226B (\ll/\gg), and the Greek/
% blackboard letters used in equations. Without an explicit \setmathfont,
% unicode-math falls back to lmroman text font which lacks several glyphs
% and emits "Missing character" warnings on every \mid in math mode.

% Slide layout defaults for warning-clean scientific decks.
\usepackage{etoolbox}
\IfFileExists{xurl.sty}{\usepackage{xurl}}{}
\IfFileExists{seqsplit.sty}{\usepackage{seqsplit}}{\newcommand{\seqsplit}[1]{#1}}
\protected\def\breakseq#1{\seqsplit{#1}}
\protected\def\breaktt#1{\begingroup\ttfamily\seqsplit{#1}\endgroup}
\setlength{\emergencystretch}{6em}
\tolerance=5000
\hbadness=10000
\hfuzz=1pt
\setlength{\tabcolsep}{2pt}
\AtBeginEnvironment{longtable}{\tiny\renewcommand{\arraystretch}{0.86}\setlength{\tabcolsep}{1pt}}
\AtBeginEnvironment{tabular}{\tiny\renewcommand{\arraystretch}{0.86}\setlength{\tabcolsep}{1pt}}
\AtBeginEnvironment{equation}{\tiny}
\AtBeginEnvironment{equation*}{\tiny}
\AtBeginEnvironment{align}{\tiny}
\AtBeginEnvironment{align*}{\tiny}
\AtBeginEnvironment{itemize}{\footnotesize}
\AtBeginEnvironment{enumerate}{\footnotesize}
\AtBeginEnvironment{description}{\footnotesize}
\setbeamerfont{caption}{size=\tiny}
\setbeamerfont{caption name}{size=\tiny}
\setbeamerfont{normal text}{size=\small}
\setbeamerfont{frametitle}{size=\small}
\setbeamerfont{section title}{size=\footnotesize}
\setbeamerfont{subsection title}{size=\footnotesize}
\setbeamertemplate{section page}{%
  \centering
  \begin{beamercolorbox}[sep=12pt,center,wd=\paperwidth]{section title}
    \parbox{0.86\paperwidth}{\centering\usebeamerfont{section title}\insertsection\par}
  \end{beamercolorbox}
}
\setbeamertemplate{subsection page}{%
  \centering
  \begin{beamercolorbox}[sep=8pt,center,wd=\paperwidth]{subsection title}
    \parbox{0.86\paperwidth}{\centering\usebeamerfont{subsection title}\insertsubsection\par}
  \end{beamercolorbox}
}
\setlength{\abovecaptionskip}{2pt}
\setlength{\belowcaptionskip}{0pt}

% Natbib and cross-reference fallbacks — slides don't load natbib
% or cleveref, but combined-PDF manuscript prose may emit these
% commands. The fallback renders citations as a bracketed key list
% and cross-references as detokenized labels so slides don't fail on
% undefined control sequences or raw underscores. \providecommand is
% a no-op if packages load later.
\providecommand{\citep}[1]{[#1]}
\providecommand{\citet}[1]{#1}
\providecommand{\citealp}[1]{#1}
\providecommand{\citeauthor}[1]{#1}
\providecommand{\citeyear}[1]{#1}
\providecommand{\cref}[1]{\texttt{\detokenize{#1}}}
\providecommand{\Cref}[1]{\texttt{\detokenize{#1}}}
\providecommand{\crefrange}[2]{\texttt{\detokenize{#1}}--\texttt{\detokenize{#2}}}
\providecommand{\Crefrange}[2]{\texttt{\detokenize{#1}}--\texttt{\detokenize{#2}}}
\providecommand{\paragraph}[1]{\textbf{#1}\ }

\newtheorem{proposition}{Proposition}
\newtheorem{hypothesis}{Hypothesis}
\newtheorem{remark}[theorem]{Remark}
\newtheorem{axiom}{Axiom}
\newtheorem{property}{Property}
\usepackage{bookmark}
\IfFileExists{xurl.sty}{\usepackage{xurl}}{} % add URL line breaks if available
\urlstyle{same}
% fallback for those not using the hyperref driver hyperxmp:
\makeatletter
\@ifundefined{xmpquote}{\newcommand{\xmpquote}[1]{#1}}{}
\makeatother
\hypersetup{
  hidelinks,
  pdfcreator={LaTeX via pandoc}}

\author{\texorpdfstring{}{}}
\date{}

\begin{document}

\section{Scope, Related Work, and Limitations}\label{sec:scope}

\begin{frame}[allowframebreaks]{Scope and epistemic boundary}
\protect\phantomsection\label{scope-and-epistemic-boundary}
DuckRabbit is a typed, reproducible stimulus-construction platform. Its
unit of work is an immutable request that yields a canonical image,
audio buffer, video sequence, or audiovisual timeline, together with
objective measurements and provenance. It is not an exhaustive ontology
of all illusions, a perceptual theory, a participant database, or a
substitute for a controlled psychophysics experiment. The live registry
is therefore a deliberately bounded catalog: it enumerates the families
for which this release has both a generator contract and a stated
evidence boundary. Appendix A gives the complete current matrix; the
absence of a family from that matrix is not a claim that the family is
unimportant or unsupported in the wider literature.

\end{frame}

\begin{frame}[allowframebreaks]{Scope and epistemic boundary}
The central distinction is between a construction and an observation. A
DuckRabbit generator can establish the dimensions, pixel values, sample
values, frequency components, frame clock, declared synchronization
offset, encoded hash, and decoded media facts of its output. It cannot,
from those facts alone, establish what a person sees, hears, counts,
localizes, groups, or reports. Every figure caption, evidence record,
and claim-ledger entry preserves this boundary. In particular, a
literature citation establishes a source-supported statement about a
stimulus family or result under that source's conditions; it does not
certify pixel- or task-identical replication by a DuckRabbit artifact.

\end{frame}

\begin{frame}[allowframebreaks]{Scope and epistemic boundary}
The historical foundation is intentionally worldwide and layered. The
package does not present a single invention narrative: it places ancient
Greek, medieval Arabic, classical Chinese, and early-modern European
sources beside nineteenth-century psychophysics and modern experimental
work. This is a scholarly orientation for separating physical
construction, interpretation, and observer evidence; it is not a claim
that the current source set exhausts the visual, auditory, or perceptual
traditions of any region.
\end{frame}

\begin{frame}[allowframebreaks]{Visual families}
\protect\phantomsection\label{visual-families}
The modern label ``visual illusion'' sits on a longer experimental
history than the package's contemporary review sources alone suggest.
Oppel's mid-century catalogue of geometrical-optical illusions is now
accessible with translation and commentary (Wade et al. 2017), and the
primary nineteenth-century records include Zöllner's account of
crossing-line distortions and Müller-Lyer's report of the arrow-wing
configuration (Zöllner 1860; Müller-Lyer 1889). These sources are
treated as historical anchors, not as evidence that a present-day raster
is identical to an original plate or that its observer effect is
invariant across conditions.

\end{frame}

\begin{frame}[allowframebreaks]{Visual families}
Gregory's classification is useful as a historical and conceptual
orientation, but it is not a universally agreed ontology. DuckRabbit
consequently stores visual modality, mechanism, perceptual signature,
stimulus requirements, and evidence status as separate facets rather
than treating a single label as an explanation (Gregory 1997). The
visual atlas covers the implemented families currently registered in the
package: ambiguous figure, contrast, apparent motion, Müller-Lyer,
Poggendorff, Ponzo, Kanizsa-type subjective contour, Ebbinghaus
contextual size, and Zöllner/Judd orientation-related geometry.
``Coverage'' here means one deterministic representative per live entry,
not a claim that one raster captures the historical stimulus space.

\end{frame}

\begin{frame}[allowframebreaks]{Visual families}
The duck/rabbit construction is an explicit example of why that
distinction matters. Brugger's historical analysis discusses variation
among figure variants and observers; DuckRabbit therefore exposes a
blend parameter and labels its output as a construction rather than
promising a fixed alternation rate or universal bistability (Brugger
1999). The Müller-Lyer entry uses typed arrow geometry. Its engineering
form is compatible with a family whose image statistics have been
analyzed as potentially informative about image-source relationships,
but the implementation does not adjudicate that account or infer a
perceived-length report (Müller-Lyer 1889; Howe and Purves 2005).

\end{frame}

\begin{frame}[allowframebreaks]{Visual families}
The Poggendorff generator makes the occluding geometry and virtual-line
orientation explicit. This is appropriate to a literature in which
orientation estimation and filtering accounts are theoretically
relevant, while avoiding the stronger claim that a particular raster
reproduces a published bias without the same observers, viewing
conditions, and response task (Morgan 1999). The Ponzo construction is
similarly bounded: the converging context and test bars are
deterministic, whereas the literature contains multiple explanations of
Ponzo-like effects rather than one settled mechanism (Fisher 1967;
Yildiz et al. 2022).

\end{frame}

\begin{frame}[allowframebreaks]{Visual families}
Kanizsa-style subjective contours are represented as an illusory-contour
construction with explicit inducer geometry. The historical account
motivates the family label, while modern reviews place contour
completion within a wider literature on grouping, figure-ground
organization, attention, and neural mechanisms (Kanizsa 1976; Wagemans
et al. 2012). Neither source licenses a claim that the generated mask
produces a uniform contour percept across observers. The Ebbinghaus
entry controls target and surround geometry; classic work demonstrates
that judged size depends on context, contour, and comparison conditions,
while later work shows that motion and other parameters can alter the
effect (Weintraub 1979; Mruczek et al. 2015). DuckRabbit's static
default is therefore an engineering baseline rather than a
task-identical replication. Zöllner/Judd geometry is likewise retained
as a source-linked line arrangement. Zöllner's 1860 report supplies the
\end{frame}

\begin{frame}[allowframebreaks]{Visual families}
historical primary anchor, while later work supplies a modern analysis
of spatial filtering (Zöllner 1860; Earle and Maskell 1995). Its typed
line lengths and orientations make the spatial stimulus auditable, while
orientation judgments remain outside the package contract.

Finally, apparent motion is included in the visual coverage panel
because its defining engineering object is temporal succession, not
merely a static pattern. The package verifies frame order, frame rate,
timestamps, and frame-to-frame differences. Wertheimer's foundational
work and Sekuler's later analysis motivate the family distinction;
neither source turns a generated frame strip into a universal report of
motion (Wertheimer 1912; Sekuler 1996).
\end{frame}

\begin{frame}[allowframebreaks]{Auditory families}
\protect\phantomsection\label{auditory-families}
The auditory catalog separates harmonic construction, spectral
completion, pitch-class context, dichotic channel assignment, and
continuity/masking. A Shepard-like signal is represented through
additive partials and explicit envelopes, following a historical
literature of tone psychology and later work on assimilation to an
internalized musical scale (Stumpf 1883; Shepard and Jordan 1984). The
canonical buffer exposes sample rate, amplitude bounds, channel count,
partial frequencies, and envelope parameters. It does not determine a
listener's perceived pitch height or direction.

\end{frame}

\begin{frame}[allowframebreaks]{Auditory families}
The missing-fundamental construction removes a low component while
retaining harmonically related partials. This makes the spectral
condition reproducible; the associated pitch interpretation remains a
listener-level question (Zatorre 2005). Tritone and octave entries
preserve channel structure, phase, frequency relationships, and timing
in typed parameters. Deutsch's foundational accounts and Repp's analysis
of spectral-envelope and context effects motivate the evidence records,
while also making listener and context dependence central limitations
(Deutsch 1986, 1974; Repp 1997).

\end{frame}

\begin{frame}[allowframebreaks]{Auditory families}
Auditory continuity is represented as an interrupted target with a typed
masker and gap. Warren's classic perceptual-restoration result and
Riecke and colleagues' separation of sensory and decisional
contributions justify the family's inclusion, but a generated masker is
not evidence that a listener will report an uninterrupted sound (Warren
1970; Riecke et al. 2011). The audio atlas therefore reports RMS, peak,
spectral centroid, bandwidth, channel layout, and exact sample-level
provenance---not continuity, pitch, or stream judgments.
\end{frame}

\begin{frame}[fragile,allowframebreaks]{Audiovisual and temporal binding
families}
\protect\phantomsection\label{audiovisual-and-temporal-binding-families}
Audiovisual constructions require a shared clock and an explicit sign
convention for offsets. The sound-induced-flash entry creates a typed
visual event and one or more audio events; the evidence record links it
to the primary demonstration and to a review of the broader literature
(Shams et al. 2000; Hirst et al. 2020). The package can verify event
timestamps, sample/frame clocks, and declared offsets. It does not infer
a reported flash count, and it does not assume that the same temporal
window applies across displays, headphones, latencies, or observers.

\end{frame}

\begin{frame}[fragile,allowframebreaks]{Audiovisual and temporal binding
families}
The spatial ventriloquist construction treats spatial discrepancy as a
parameter and records the channel and display requirements. Reviews
describe the ventriloquist illusion as a tool for studying multisensory
processing, but the review's scope is not a license to claim spatial
capture from a file alone (Bruns 2019). More general audiovisual
accounts describe integration and segregation as a causal-inference
problem shaped by temporal regularities and signal reliability (Noppeney
and Lee 2018). That framework clarifies why a declared offset is an
experimental input, not an observer-level outcome. Temporal
ventriloquism receives a separate temporal-binding signature rather than
being collapsed into spatial capture. Vroomen and de Gelder manipulated
sound--flash timing in a flash-lag task and reported timing-dependent
changes under those experimental conditions; Hartcher-O'Brien and Alais
\end{frame}

\begin{frame}[fragile,allowframebreaks]{Audiovisual and temporal binding
families}
studied temporal ventriloquism in a purely temporal context (Vroomen and
Gelder 2004; Hartcher-O'Brien and Alais 2011). DuckRabbit implements the
stimulus-side timing contract and leaves the observer-side
temporal-recalibration estimate to a future study.

McGurk remains \texttt{input\_required}. The classic speech study
motivates the family, but a lawful and reproducible implementation
requires checksummed speech and video fixtures, licensing or consent
records, a precise preprocessing contract, and an ethical validation
protocol (McGurk and MacDonald 1976). Cataloguing the gap is more
informative than silently substituting an unrelated synthetic voice or
claiming that a generic audiovisual mismatch is a McGurk replication.
\end{frame}

\begin{frame}[allowframebreaks]{From literature to engineering contract}
\protect\phantomsection\label{from-literature-to-engineering-contract}
For each entry, the evidence matrix records a source role, exact
source-supported claim, engineering basis, limitation, and audit status.
The roles distinguish primary demonstration, review or synthesis,
theoretical account, engineering basis, limitation, and input gap. This
prevents three common category errors: treating a theory as settled
mechanism, treating a review as validation of new code, and treating a
generator's deterministic output as a participant result.

\end{frame}

\begin{frame}[allowframebreaks]{From literature to engineering contract}
The package's literature fidelity is therefore best described as
``family-linked, parameterized engineering construction.'' Some defaults
are historically motivated; none should be read as a claim of pixel
identity unless that identity is separately established. The checked-in
scholarship snapshot provides offline structural validation, including
citation-key and DOI/URL format checks. The explicit network audit adds
resolver observations and metadata matching; reachability alone is not
treated as bibliographic verification. This is a reproducible evidence
boundary, not a claim that every source is equally accessible or that
theoretical disputes have been resolved.
\end{frame}

\begin{frame}[fragile,allowframebreaks]{Limitations and future observer
work}
\protect\phantomsection\label{limitations-and-future-observer-work}
The package does not claim clinical validity, universal effect sizes,
cross-cultural invariance, perceptual equivalence across displays or
headphones, or observer-level truth from objective media metrics.
Playback level, gamma and display calibration, viewing distance, refresh
rate, stereo separation, room acoustics, audio transducer response,
attention, expectation, language, expertise, and task can all matter.
Codec behavior and device latency can introduce additional differences
even when decoded media facts match within the declared tolerance.

\end{frame}

\begin{frame}[fragile,allowframebreaks]{Limitations and future observer
work}
The synthetic observer diagnostic is intentionally not a substitute for
a real vision model or human data. It is a serialized, hand-specified
feature function with \breaktt{training\_data=none},
\breaktt{calibration\ =\ analytic}, and \breaktt{human\_data=false}; its
metamorphic and sensitivity checks test orchestration, not visual
consciousness or human discrimination. A future observer study must
pre-register the response scale, reference condition, estimand,
exclusions, missing-data rule, randomization seed, display and playback
controls, and uncertainty procedure. It must also report the participant
and item sampling frame rather than importing a model-output probability
as an assumed effect size. The observer harness is consequently
study-ready scaffolding, not a result set.

\end{frame}

\begin{frame}[fragile,allowframebreaks]{Limitations and future observer
work}
No participant data are bundled with DuckRabbit. The data-availability
boundary is intentional: canonical artifacts, encoded fixtures where
lawful, source-data sidecars, and deterministic synthetic diagnostics
are software outputs; observer outcomes require separately governed data
collection. The package is authored by Daniel Ari Friedman and is
archived under DOI \breaktt{10.5281/zenodo.21419693}; release metadata
and software-citation guidance are generated from the same identity
contract as the manuscript.
\end{frame}

\end{document}
